\section{Resultados}

\subsection{Matriz de riesgo consolidada}

\begin{table}[htbp]
\centering
\caption{Resumen de escenarios de fuego dominantes (HAZOP + What-If)}
\label{tab:matriz_riesgo}
\setcounter{itemcount}{0}
\begin{tabularx}{\textwidth}{@{}NXcccl@{}}
\toprule
\multicolumn{1}{c}{\textbf{Item}} & \textbf{Escenario} & \textbf{Sev.} & \textbf{Prob.} & \textbf{Riesgo} & \textbf{Salvaguarda propuesta} \\
\midrule
& Incendio en patio de bagazo & 4 & 3 & 12 (Alto) & Monitores + rociadores perimetrales \\
& Deflagración en silo de bagazo & 5 & 2 & 10 (Alto) & Venteo NFPA 68 + inertización \\
& Fuego en aceite de turbina & 4 & 3 & 12 (Alto) & CO\textsubscript{2} + agua pulverizada \\
& Cortocircuito en transformador & 4 & 2 & 8 (Moderado) & Diluvio NFPA 15 \\
& Incendio en sala de control & 3 & 2 & 6 (Moderado) & Agente limpio NFPA 2001 \\
\bottomrule
\end{tabularx}
\end{table}

\subsection{Cálculos hidráulicos — red de rociadores}

\begin{table}[htbp]
\centering
\caption{Resultados del nodal hidráulico del área más remota}
\label{tab:hidraulica_resultados}
\setcounter{itemcount}{0}
\begin{tabularx}{\textwidth}{@{}NXlcl@{}}
\toprule
\multicolumn{1}{c}{\textbf{Item}} & \textbf{Variable} & \textbf{Valor} & \textbf{Unidad} & \textbf{Referencia} \\
\midrule
& Caudal total demanda & 500 & gpm & NFPA 13 \\
& Presión en la conexión & 85 & psi & \\
& $\Delta P$ red + elevación & 42 & psi & Ec.\ \ref{eq:hw} \\
& Rociadores operando & 15 & -- & Área más remota \\
& Margen de seguridad & 12 & \% & \\
\bottomrule
\end{tabularx}
\end{table}

\subsection{Dimensionamiento de bomba (NFPA 20)}

Se selecciona una bomba centrífuga horizontal de partición axial con capacidad nominal de 500 gpm a 100 psi. Se verifica que el punto de \textit{churn} no exceda 140\,\% de la presión de diseño y que al 150\,\% del caudal la presión sea $\ge$ 65\,\% del diseño.

\subsection{Detección y alarma (NFPA 72)}

\begin{itemize}
\item Capacidad de batería calculada: $C_{\text{Ah}} = (I_{sb}\,t_{sb} + I_{al}\,t_{al}) \cdot 1{,}20$.
\item Caída de tensión máxima en NAC: $<$ 10\,\% al dispositivo más alejado.
\item Nivel sonoro en peor punto: $\ge$ 15 dB sobre ruido ambiente.
\end{itemize}

\subsection{Agente limpio (NFPA 2001)}

Masa de agente \agenteLimpio{} para el recinto protegido con concentración de diseño \concentracionDiseno{}:

\begin{equation*}
m = \frac{V}{k_1 + k_2 T} \cdot \frac{C}{100 - C}
\end{equation*}

Verificación: $C_{\text{diseño}} \le \text{NOAEL}$ para ocupación normal.

\subsection{Clasificación de áreas NEC}

El resumen por área se encuentra en el Anexo correspondiente; ver también \texttt{data/clasificacion\_areas.csv}.
